\documentclass[11pt]{scrartcl}
\usepackage[a4paper,margin=2.6cm]{geometry}
\usepackage{../common/rlstudy}
\title{Extended reading --- Week 8: A2C, GAE, and TRPO}
\date{\today}
\begin{document}
\maketitle

\section*{How to use these readings}
Start with the teaching guide's three-step GAE calculation and two-dimensional
Fisher solve. For each paper below, identify the mathematical objective,
the implemented approximation and the evaluation evidence separately.
The reading questions are course exercises, not quotations from the papers.

\section*{1. Mnih et al. (2016): Asynchronous Methods for Deep Reinforcement Learning}
\textbf{Source:} ICML 2016, \url{https://arxiv.org/abs/1602.01783}.

\paragraph{Contribution.} The paper studies asynchronous deep-RL algorithms,
including A3C, with parallel actor-learners. Parallel collection can diversify
experience without a replay buffer. The synchronous A2C used in this course
shares the actor-critic idea but waits to combine a batch of rollouts.

\paragraph{Read critically.} Independent environment instances do not make
successive transitions within a trajectory independent. Identify which
parameters each actor uses and when those parameters change. A PPO collector
can resemble an A2C collector while using a different optimisation objective
and repeated minibatch passes.

\paragraph{Notebook connection.} Inspect the time and environment axes before
flattening a rollout. Explain why residuals from two different environments
must not be joined into one GAE trace. Draw eight separate short trajectories,
not one long trajectory with eight times as many transitions.

\section*{2. Schulman et al. (2015/2016): Generalized Advantage Estimation}
\textbf{Source:} ICLR 2016, \url{https://arxiv.org/abs/1506.02438}.

\paragraph{Contribution.} GAE combines temporal-difference residuals with
exponentially decreasing weights. The paper studies the bias and variance
introduced by using values and discounting in policy-gradient estimation.
The course's closed-form derivation changes the order of summation; it does
not independently establish an unbiased full training procedure.

\paragraph{Read critically.} Separate the return objective's discount from
lambda's estimator choice. A learned critic can be inaccurate, and a finite
rollout can retain an endpoint bootstrap. Neither lambda zero nor lambda one
has a universally best mean-squared error. Empirical settings need their
measurement context.

\paragraph{Notebook connection.} For residuals $(1,2)$ and $\gamma\lambda=0.72$,
compute the first coefficient sum $1+0.72(2)=2.44$. Now change only the boundary:
which residual changes if a final value is kept instead of terminal zero?
Compare the reconstructed return with the direct finite reward sum.

\section*{3. Kakade (2001): A Natural Policy Gradient}
\textbf{Source:} NeurIPS 2001,
\url{https://proceedings.neurips.cc/paper/2001/hash/4b86abe48d358ecf194c56c69108433e-Abstract.html}.

\paragraph{Contribution.} Natural policy gradients account for the local
geometry of policy distributions through a Fisher metric. Under appropriate
conditions, the direction solves $Fx=g$. Compatible function approximation
connects a particular advantage regression with that direction.

\paragraph{Read critically.} State which policy and state weighting define
$F$. A singular empirical matrix needs additional treatment. A connection
between KL geometry and natural gradients does not make DPO a Fisher-inverse
approximation or derive its beta setting; DPO has its own preference objective.

\paragraph{Notebook connection.} With $F=\operatorname{diag}(4,1)$ and
$g=(2,1)$, solve $x=(0.5,1)$. Under a quadratic KL budget $0.01$, scale it to
$u=(0.05,0.1)$. Check both $g^Tu=0.2$ and $u^TFu/2=0.01$.
A small displacement in a sensitive coordinate can produce the same local
policy change as a larger displacement elsewhere.

\section*{4. Schulman et al. (2015): Trust Region Policy Optimization}
\textbf{Source:} ICML 2015, \url{https://arxiv.org/abs/1502.05477}.

\paragraph{Contribution.} The paper develops a policy-improvement bound and
then makes approximations to obtain practical TRPO. The theoretical
maximum-state divergence and exact advantages differ from the sampled
average-KL constraint and estimated advantages used in implementation.

\paragraph{Read critically.} Backtracking checks a sampled surrogate and
sampled KL; it does not remove all approximation or sampling error.
Do not read a theorem about the ideal bound as an unconditional guarantee
for a neural implementation. PPO clipping is another objective design,
not an exact solution of the TRPO constraint.

\paragraph{Notebook connection.} Reuse the guide's Fisher solve. Perturb the
step size and compute the local quadratic KL, then distinguish that number
from KL measured using the actual changed policy. Explain what a rare state
with a large action-distribution change can hide in a small average.

\section*{5. Engstrom et al. (2020): Implementation Matters in Deep Policy Gradients}
\textbf{Source:} ICLR 2020, \url{https://arxiv.org/abs/2005.12729}.

\paragraph{Contribution.} The experiments show that implementation choices
can substantially affect PPO/TRPO comparisons and the attribution of gains.
The result is evidence for the tested settings, not a universal percentage
of improvement attributable to implementation across every task.

\paragraph{Read critically.} Record the interaction budget, evaluation policy,
normalisation choices and random seeds for an ablation. Changing several
choices together cannot isolate the effect of one component. Similar
collector code also does not make two update objectives equivalent.

\paragraph{Notebook connection.} Before a costly lambda sweep, reproduce the
small deterministic calculation. Then specify which outcome a training sweep
would test: evaluation return, variance of gradient estimates, or critic error.
These require different measurements. A plot of scalar losses is not a
measurement of parameter-gradient variance.

\section*{After reading}
Write one paragraph connecting collection, advantage estimation and policy
update. Name one quantity observed from the environment, one predicted by a
critic, and one checked by a line search. State which claims are algebra,
which rely on assumptions, and which require new experiments.
\end{document}
